\problemname{Trafikverkets plan}
% TODO
Under vintersårstiden sker det på tok för många krockar i lönnköping.
Staden är ett träd. Rikta varje kant. Kolla att alla kan ta sig till jobbet.
På kvällen riktar de i motsatt håll, allt blir perfekt. Gods plan referens? 

\section*{Indata}
Den första raden i indatan innehåller heltalet $N$ ($1 \le N \le 10^5$), längden på protein $T$.

Nästa rad innehåller en sträng av längd $N$, protein $T$.

På raden därefter följer $M$ ($1 \le M \le 10^6$), antalet protein Björn kan tillverka.

Därefter följer $M$ rader, som vardera innehåller en sträng. Denna sträng representerar ett protein han
kan tillverka. Det är garanterat att den totala längden av de $M$ strängarna inte överstiger $10^6$.

Alla strängar i indatan består av bokstäver \texttt{a}-\texttt{z}. Det är garanterat att Björn kan konstruera $T$ med hjälp av de $M$ strängar som Björn kan tillverka.

\section*{Utdata}
Skriv ut ett heltal: det minsta antalet protein som Björn behöver producera för att skapa protein $T$.


\section*{Poängsättning}
Din lösning kommer att testas på en mängd testfallsgrupper.
För att få poäng för en grupp så måste du klara alla testfall i gruppen.

\noindent
\begin{tabular}{| l | l | p{12cm} |}
  \hline
  \textbf{Grupp} & \textbf{Poäng} & \textbf{Gränser} \\ \hline
  $1$    & $10$       & $N \leq 20, M \leq 5$ \\ \hline
  $2$    & $20$       & $N, M \leq 100$ \\ \hline
  $3$    & $20$       & $M \leq 100$ \\ \hline
  $4$    & $40$       & $N \leq 10,000$ \\ \hline
  $5$    & $10$       & Inga ytterligare begränsningar. \\ \hline
\end{tabular}

\section*{Förklaring av exempelfall 1}
Genom att skapa proteinet \texttt{abc} två gånger kan han sätta ihop de och skapa protein $T$.
